% Generated draft: requires mathematical and source review.
\documentclass[11pt]{article}
\usepackage[margin=25mm]{geometry}
\usepackage{fontspec}
\setmainfont{Latin Modern Roman}
\usepackage{amsmath,amssymb,mathtools,unicode-math}
\setmathfont{Latin Modern Math}
\usepackage{enumitem,needspace,xurl,hyperref}
\setlength{\parindent}{0pt}
\setlength{\parskip}{5pt}
\setlength{\emergencystretch}{3em}
\newcommand{\N}{\mathbb{N}}
\newcommand{\Z}{\mathbb{Z}}
\newcommand{\Q}{\mathbb{Q}}
\newcommand{\R}{\mathbb{R}}
\newcommand{\C}{\mathbb{C}}
\newcommand{\F}{\mathbb{F}}
\newcommand{\A}{\mathbb{A}}
\newcommand{\PP}{\mathbb{P}}
\newcommand{\E}{\mathbb{E}}
\newcommand{\eps}{\varepsilon}
\newcommand{\dd}{\,\mathrm{d}}
\begin{document}
\section*{UnsolvedMath ID: 20001299}
\textbf{Source problem code:} catalogue entry
\subsection{Definitions and mathematical statement}
A full-dimensional origin-symmetric zonoid is a Hausdorff limit of finite Minkowski sums of centered line segments in \(\R^n\). For an origin-symmetric convex body \(Z\), let \(Z^\circ=\{y:x\cdot y\le1\text{ for all }x\in Z\}\), and let \(B_2^n\) be the Euclidean unit ball. Use the Banach--Mazur convention
\[
d_{\rm BM}(K,L)=\inf\{\lambda\ge1:\ K\subseteq TL\subseteq\lambda K\text{ for some invertible linear }T\}.
\]
R. Schneider's question asks whether, for every sequence \((Z_n)\) with both \(Z_n\subset\R^n\) and \(Z_n^\circ\) zonoids,
\[
d_{\rm BM}(Z_n,B_2^n)\longrightarrow1\quad(n\to\infty).
\]
Equivalently, does the supremum of these distances over the polar-zonoid class in dimension \(n\) tend to 1? Full dimensionality and symmetry about the origin are the standard conventions needed for a bounded polar and this distance. The source's damaged symbols \(Z_n\), \(Z_n^\circ\), \(B_2^n\) are interpreted accordingly.

The supplied aid reports the unrestricted asymptotic problem open in the literature it checked. Its special case fixes an integer \(m\ge1\), takes \(n\ge2m\), and defines the even spherical harmonic
\[
Y_{m,n}(u)=(2m)^m\prod_{j=1}^{2m}u_j,\qquad u\in S^{n-1}.
\]
Its range is \([-1,1]\). For the candidate support function
\[
h_\varepsilon(u)=1+\varepsilon Y_{m,n}(u),
\]
the aid asserts the exact criterion
\[
h_\varepsilon\text{ is the support function of a zonoid}
\quad\Longleftrightarrow\quad
|\varepsilon|\le a_{m,n},\qquad
a_{m,n}=\frac{(2m-3)!!}{\prod_{j=1}^m(n+2j-1)}.
\]
Here \((-1)!!=1\) for \(m=1\); for \(m\ge2\), the numerator is the product of positive odd integers through \(2m-3\). This criterion concerns existence as a zonoid support function, not an assertion that every \(\varepsilon\) outside the interval defines a convex body of a different class.

For each fixed pair \((m,n)\), the aid also asserts existence of \(\delta_{m,n}>0\) such that the bodies with support function \(h_\varepsilon\), for sufficiently small \(|\varepsilon|<\delta_{m,n}\), have zonoid polars as well. One may decrease this neighborhood to have \(\delta_{m,n}\le a_{m,n}\), so both body and polar are zonoids there. The summary attributes the conclusion to smooth polarity and continuity of the inverse cosine transform, but provides no explicit quantitative lower bound on \(\delta_{m,n}\) or uniform neighborhood in \(n\).

Every member of this fixed-harmonic family satisfying both zonoid conditions has the reported estimate
\[
d_{\rm BM}(Z_{n,\varepsilon},B_2^n)
\le\frac{1+a_{m,n}}{1-a_{m,n}}
=1+O_m(n^{-m})\quad(n\to\infty,\ m\text{ fixed}).
\]
The denominator is positive for the stated range of \(m,n\). The implicit constant may depend on \(m\). This proves asymptotic closeness only for the displayed one-harmonic subfamilies; it does not supply a uniform estimate for all polar-zonoid bodies or for harmonic degree growing arbitrarily with dimension.

The registry labels the entry partially solved; attempt 1 is a partial result, with problem open at the run, internal rigor label proved and novelty confidence low. Those labels accompany a machine-generated aid that is not peer reviewed and requires independent verification before citation.
\subsection{Short English statement}
Schneider asks whether a zonoid whose polar is also a zonoid must become uniformly close to an ellipsoid in high dimension. The aid proves this only for a specified fixed-degree harmonic perturbation of a ball, with an exact zonoid threshold and an unquantified local interval for polarity. The full asymptotic question remains open according to the supplied checked-literature report.
\subsection{Sources}
\begin{itemize}
\item Local UnsolvedMath registry, permanent ID \(20001299\), source code \texttt{catalogue entry}; \texttt{aim-convex-geometry-notes.json}, \texttt{notes[30]}.
\item AIM, \emph{Fourier analytic methods in convex geometry}, Problem 1 (R. Schneider): \url{https://aimath.org/WWN/fourierconvex/fourierconvex.pdf}. The supplied verification notes damaged subscript and superscript extraction; no independent PDF visit is claimed.
\item Supplied research attempt 1 and registry update August 14, 2026. The original literature field is empty, and no separate authoritative publication reference for the exact one-harmonic criterion or polarity neighborhood is supplied.
\end{itemize}
\end{document}
